% ==== Responsável: Letícia Storch ====
\chapter{Processo de \textit{boot}}\label{cap:boot}

O \textit{boot} (ou inicialização) é o conjunto de etapas que leva o computador do estado
desligado até um sistema operacional pronto para uso. Como a memória principal é volátil, o
\textit{hardware} precisa localizar, em algum dispositivo de armazenamento, o código capaz de
carregar o núcleo (\textit{kernel}) na memória e entregar-lhe o controle \cite{tanenbaum2016}.

\section{Etapas da inicialização}\label{sec:etapas-boot}

Ao ligar a máquina, o processador começa a executar um código gravado em memória não volátil:
o \textbf{\textit{firmware}}, que pode ser a BIOS (modo legado) ou a
UEFI\footnote{Especificações UEFI. Disponível em: \url{https://uefi.org/specifications}. Acesso
em: 19 nov. 2026.}. O \textit{firmware} executa o POST (\textit{Power-On Self Test}), que testa
componentes básicos como memória e teclado, e em seguida procura um dispositivo de
\textit{boot}. A \autoref{fig:boot} resume o fluxo e a \autoref{tab:bios-uefi} compara os dois
tipos de \textit{firmware}.

\begin{figure}[H]
    \centering
    \caption{Fluxo simplificado do processo de \textit{boot} no Linux}
    \label{fig:boot}
    \begin{tikzpicture}[node distance=6mm and 8mm,
        box/.style={rectangle, draw, rounded corners, align=center, fill=gray!10,
                    minimum width=2.6cm, minimum height=1cm, font=\small},
        arr/.style={->, >=stealth, thick}]
        \node[box] (fw) {\textit{Firmware}\\BIOS/UEFI};
        \node[box, right=of fw] (post) {POST};
        \node[box, right=of post] (bl) {\textit{Bootloader}\\(GRUB)};
        \node[box, below=of bl] (k) {\textit{Kernel}\\+ \textit{initramfs}};
        \node[box, left=of k] (init) {\textit{Init}\\(\texttt{systemd})};
        \node[box, left=of init] (user) {Espaço do\\usuário};
        \draw[arr] (fw) -- (post);
        \draw[arr] (post) -- (bl);
        \draw[arr] (bl) -- (k);
        \draw[arr] (k) -- (init);
        \draw[arr] (init) -- (user);
    \end{tikzpicture}
    \legend{Fonte: os autores (2026).}
\end{figure}

\begin{table}[H]
    \centering
    \caption{Comparação entre BIOS e UEFI}
    \label{tab:bios-uefi}
    \begin{tabularx}{\linewidth}{lXX}
        \toprule
        \textbf{Característica} & \textbf{BIOS (legado)} & \textbf{UEFI} \\
        \midrule
        Esquema de partição & MBR & GPT (também aceita MBR) \\
        Tamanho máximo de disco & Cerca de 2\,TiB & Acima de 2\,TiB \\
        Local do \textit{bootloader} & Primeiro setor do disco (MBR, 512 bytes) & Partição de sistema EFI (ESP, em FAT) \\
        Interface de configuração & Texto & Texto ou gráfica \\
        \textit{Secure Boot} & Não & Sim \\
        \bottomrule
    \end{tabularx}
    \legend{Fonte: os autores (2026), com base em \citeonline{silberschatz2015}.}
\end{table}

\section{\textit{Bootloader}}\label{sec:bootloader}

O \textit{bootloader} é o programa responsável por carregar o \textit{kernel}. No Linux, o mais
comum é o GRUB\footnote{Manual do GRUB. Disponível em:
\url{https://www.gnu.org/software/grub/manual/grub/}. Acesso em: 19 nov. 2026.}, que permite
escolher entre os sistemas instalados e passar parâmetros ao \textit{kernel}. Ele carrega também
o \textit{initramfs}, um sistema de arquivos temporário na memória com os \textit{drivers}
necessários para montar o sistema de arquivos raiz.

O menu do GRUB também dá acesso ao \textbf{modo de recuperação} (\textit{recovery mode}) do
Ubuntu, que inicia o sistema com o mínimo de serviços e oferece um terminal de administrador.
Nesse modo, o sistema de arquivos raiz é montado apenas para leitura e precisa ser remontado
com \texttt{mount -o remount,rw /} antes de qualquer alteração, como a troca de uma senha
esquecida com \texttt{passwd}.

\section{\textit{Kernel} e processo \texttt{init}}\label{sec:init}

Com o \textit{kernel} na memória, ele inicializa os dispositivos, monta o sistema de arquivos
raiz e cria o primeiro processo do espaço do usuário (PID~1). Nas distribuições modernas, esse
processo é o \texttt{systemd}, que inicia os serviços em paralelo e define o estado final do
sistema (por exemplo, modo gráfico ou modo texto). Os principais passos, em ordem, são:

\begin{enumerate}
    \item o \textit{firmware} executa o POST e localiza o dispositivo de \textit{boot};
    \item o \textit{bootloader} é carregado e exibe o menu de opções;
    \item o \textit{kernel} e o \textit{initramfs} são carregados na memória;
    \item o \textit{kernel} monta o sistema de arquivos raiz;
    \item o \texttt{init}/\texttt{systemd} inicia os serviços e o ambiente do usuário.
\end{enumerate}

Depois da inicialização, o comando \texttt{systemd-analyze} mostra quanto tempo cada fase do
\textit{boot} levou: o \textit{kernel} e o espaço do usuário e, em máquinas UEFI, também o
\textit{firmware} e o \textit{bootloader}.
